\documentclass[10pt, aspectratio=169]{beamer}
\usepackage[utf8]{inputenc}
\usepackage[T1]{fontenc}
\usepackage[french]{babel}
\usepackage{amsmath, amssymb}
\usepackage{graphicx}
\usepackage{xcolor}
\usepackage{booktabs}
\usepackage{multirow}
\usepackage{hyperref}
\usepackage{caption}
\usepackage{subcaption}
\usepackage{sourcecodepro}

% --- Style personnalisé (rouge, noir, bleu, blanc) ---
\definecolor{mainred}{RGB}{200, 0, 0}      % rouge vif
\definecolor{mainblue}{RGB}{0, 100, 200}   % bleu profond
\definecolor{mainblack}{RGB}{0, 0, 0}      % noir
\definecolor{mainwhite}{RGB}{255, 255, 255}% blanc

\setbeamercolor{title}{fg=mainwhite, bg=mainred}
\setbeamercolor{frametitle}{fg=mainwhite, bg=mainblue}
\setbeamercolor{structure}{fg=mainred}
\setbeamercolor{normal text}{fg=mainblack, bg=mainwhite}
\setbeamercolor{alerted text}{fg=mainred}
\setbeamercolor{block title}{fg=mainwhite, bg=mainblue}
\setbeamercolor{block body}{fg=mainblack, bg=mainwhite!90!gray}

\setbeamertemplate{navigation symbols}{}
\setbeamertemplate{footline}[frame number]
\setbeamertemplate{blocks}[rounded][shadow=true]

\title{\textcolor{mainwhite}{Analyse de Sentiments de Critiques de Films}}
\subtitle{Modèle LSTM bidirectionnel + embeddings GloVe}
\author{Awa FAYE}
\institute{Université Iba Der THIAM de Thiès \\ UFR SES / SET - Master SDA (Option IA)}
\date{Année académique 2025-2026}

\begin{document}

\begin{frame}
\titlepage
\end{frame}

\begin{frame}{Plan}
\tableofcontents
\end{frame}

\section{Introduction et objectifs}
\begin{frame}{Contexte}
\begin{itemize}
    \item Classification automatique de critiques de films (IMDb) : sentiment \textbf{positif} ou \textbf{négatif}.
    \item Objectif : dépasser une simple démonstration pour obtenir un classifieur \textbf{robuste} et \textbf{prêt à la production}.
    \item Utilisation de techniques avancées : \textcolor{mainblue}{LSTM bidirectionnel}, \textcolor{mainred}{embeddings pré-entraînés GloVe}, \textcolor{mainblue}{fine-tuning}.
\end{itemize}
\end{frame}

\section{Données et prétraitement}
\begin{frame}{Le jeu de données IMDb}
\begin{columns}[T]
\column{0.5\textwidth}
\begin{itemize}
    \item 50\,000 critiques de films (25\,000 train / 25\,000 test).
    \item Équilibre parfait : 50\% positifs, 50\% négatifs.
    \item Distribution des longueurs (mots) – \textcolor{mainblue}{90e percentile} = 450 mots.
\end{itemize}
\column{0.5\textwidth}
\centering
\includegraphics[width=\linewidth]{figures/distribution_longueurs.png}
\captionof{figure}{Histogramme des longueurs des critiques}
\end{columns}
\end{frame}

\begin{frame}{Nettoyage et tokenisation}
\begin{itemize}
    \item Suppression des balises HTML (ex. \texttt{<br />}).
    \item Suppression de la ponctuation, des chiffres, mise en minuscules.
    \item Tokenisation avec \texttt{Tokenizer} (Keras) – \textcolor{mainred}{10\,000 mots les plus fréquents}.
    \item Séquence padding à \texttt{maxlen = 450}.
    \item Division : 80\% entraînement (20\,000), 20\% validation (5\,000).
\end{itemize}
\end{frame}

\section{Architecture du modèle}
\begin{frame}{Modèle LSTM bidirectionnel avec GloVe}
\begin{block}{Couches principales}
\begin{itemize}
    \item \textbf{Embedding} : matrice GloVe (100 dimensions) – gelée au début.
    \item \textbf{Bidirectional LSTM} : 64 unités, dropout 0.2.
    \item \textbf{Dropout} : 0.5.
    \item \textbf{Dense} : 32 neurones (ReLU).
    \item \textbf{Sortie} : sigmoïde (binaire).
\end{itemize}
\end{block}
\vspace{0.2cm}
\centering
\textcolor{mainblue}{Total params : 1\,320\,065} \quad | \quad \textcolor{mainred}{Trainable après fine-tuning : 320\,065}
\end{frame}

\begin{frame}{Fine-tuning progressif}
\begin{enumerate}
    \item \textbf{Phase 1} (embeddings gelés) : 20 époques, learning rate 0.001.
    \item \textbf{Phase 2} (fine-tuning) : embeddings dégelés, 5 époques, learning rate réduit.
    \item Callbacks : \texttt{EarlyStopping} (patience=3) et \texttt{ReduceLROnPlateau}.
\end{enumerate}
\end{frame}

\section{Résultats expérimentaux}
\begin{frame}{Courbes d’apprentissage}
\begin{columns}[T]
\column{0.5\textwidth}
\centering
\includegraphics[width=\linewidth]{figures/learning_curves_phase1.png}
\captionof{figure}{Phase 1 – Embeddings gelés}
\column{0.5\textwidth}
\centering
\includegraphics[width=\linewidth]{figures/learning_curves_finetune.png}
\captionof{figure}{Phase 2 – Fine-tuning}
\end{columns}
\end{frame}

\begin{frame}{Évaluation sur l’ensemble de test}
\begin{itemize}
    \item \textbf{Accuracy globale} : \textcolor{mainred}{89,87\%}.
    \item Matrice de confusion :
\end{itemize}
\centering
\includegraphics[width=0.6\linewidth]{figures/confusion_matrix.png}
\captionof{figure}{Matrice de confusion – test (25\,000 critiques)}
\end{frame}

\begin{frame}{Rapport de classification détaillé}
\centering
\small
\begin{tabular}{lcccc}
\toprule
Classe & Précision & Rappel & F1-score & Support \\
\midrule
Négatif & 0.91 & 0.88 & 0.90 & 12\,500 \\
Positif & 0.89 & 0.92 & 0.90 & 12\,500 \\
\midrule
Accuracy & \multicolumn{4}{c}{0.90} \\
Macro avg & 0.90 & 0.90 & 0.90 & 25\,000 \\
\bottomrule
\end{tabular}
\vspace{0.5cm}

\alert{Remarque} : Performance équilibrée entre les deux classes – pas de biais significatif.
\end{frame}

\section{Exemple de prédiction}
\begin{frame}{Test sur une critique réelle}
\begin{block}{Critique (extrait)}
\small
\textit{“This is one of the first films I can remember... very touching and I gotta go.”}
\end{block}
\vspace{0.2cm}
\textbf{Résultat} : \\
\centering
\colorbox{mainblue!20}{\parbox{0.8\linewidth}{\centering
Probabilité de positivité : \textcolor{mainred}{0.9899} \\
\large \textcolor{mainblue}{Sentiment : POSITIF}
}}
\vspace{0.3cm}
\flushleft
Le modèle capture correctement la tonalité globalement positive malgré quelques expressions mitigées.
\end{frame}

\section{Conclusion et perspectives}
\begin{frame}{Conclusion}
\begin{itemize}
    \item Un classifieur \textbf{robuste} (89,9\% de test accuracy) utilisant un \textcolor{mainblue}{LSTM bidirectionnel} et des \textcolor{mainred}{embeddings GloVe}.
    \item Le fine-tuning des embeddings améliore légèrement la généralisation (validation \textasciitilde 90\%).
    \item Modèle et tokenizer sauvegardés (\texttt{.h5} et \texttt{.pkl}) – prêts pour une \textbf{application web} (Streamlit).
\end{itemize}
\vspace{0.5cm}
\begin{block}{Perspectives}
\begin{itemize}
    \item Essayer des embeddings contextuels (BERT, CamemBERT).
    \item Déploiement sur un cloud avec API REST.
    \item Analyse de sentiments en français.
\end{itemize}
\end{block}
\end{frame}

\begin{frame}{Merci de votre attention}
\centering
\Huge \textcolor{mainred}{Questions ?}
\vspace{1cm}

\footnotesize
\textbf{Références} : \\
\url{http://ai.stanford.edu/~amaas/data/sentiment/} \\
\url{https://nlp.stanford.edu/projects/glove/} \\
Code source : GitHub – \texttt{imdb\_sentiment\_analysis}
\end{frame}

\end{document}